\documentclass[10pt,a4paper]{amsart}
\usepackage[margin=19mm]{geometry}
\usepackage{amsmath,amssymb,mathtools}
\usepackage[scr=rsfs,cal=euler]{mathalfa}
\usepackage{xfrac}
\usepackage[unicode,hidelinks,bookmarks=false]{hyperref}
\newcommand{\ie}{i.e.\ }
\newcommand{\cf}{cf.\ }
\newcommand{\resp}{respectively\ }
\newcommand{\Subsection}{\subsection}
\providecommand{\nobreakdash}{\nobreak\hbox{-}}
\setlength{\emergencystretch}{2em}
\multlinegap0pt
\usepackage{bm}
\usepackage{bbm}
\usepackage{dsfont}
\usepackage{xspace}
\usepackage{cancel}
\usepackage{mathtools}
\usepackage{amsthm}
\makeatletter
\@ifundefined{theorem}{\newtheorem{theorem}{Theorem}}{}
\@ifundefined{proposition}{\newtheorem{proposition}[theorem]{Proposition}}{}
\makeatother
\newcommand{\FDP}{\operatorname{FDP}}
\newcommand{\F}{\mathcal F}
\newcommand{\Rfam}{\mathfrak R}
\newcommand{\one}{\mathbf 1}
\title[Dynamic $e$-closure for online hypotheses with any-time-vali]{Dynamic $e$-closure for online hypotheses with any-time-valid evidence: closure principles and projective mergers}
\author{Rianne de Heide}
\date{Source: arXiv:2608.09927v1 (CC BY 4.0)}
\begin{document}
\maketitle

\noindent\emph{Source and licence.} Dynamic $e$-closure for online hypotheses with any-time-valid evidence: closure principles and projective mergers. Version arXiv:2608.09927v1 (CC BY 4.0), distributed under the Creative Commons Attribution 4.0 International licence, as recorded in the arXiv licence metadata for this exact version and shown on the abstract page https://arxiv.org/abs/2608.09927v1. Full rights notice: licenses/arxiv-2608.09927-CC-BY-4.0.txt.\par\smallskip

\noindent\textit{Editorial scope.} Counted unit: the Proposition whose source label is \texttt{prop:coherence-necessary} in the article (source0.tex lines 391--393, source label \texttt{prop:coherence-necessary}), delivered together with the article's own complete original proof of it (source0.tex lines 395--431, one \texttt{proof} environment of 37 lines). No mathematics is written, repaired or re-derived: statement and proof are the article's own text line for line. Required context retained uncounted: none. Every definition, display and label of those lines is retained unchanged. Omitted from this package and not redistributed: 1--390, 394--394, 432--1694. Nothing inside a retained excerpt is omitted or altered: every retained body is delivered exactly as the article's lines are written, so no omission receipt of any kind accompanies this package. Citations: the retained excerpts carry no citation command at all, so no bibliography entry of the article is transcribed and no reference is redistributed. Reference adaptation: none was needed and none was made; every reference inside the delivered unit resolves inside it, so no unresolved reference or citation is produced. Printed slips kept literal: no printed word, symbol, hypothesis or number of the counted statement or of its proof was altered. Layout and class: the article's own class is \texttt{article} with packages including fontenc, geometry, amsmath, amssymb, amsthm, mathtools, natbib, mathrsfs, bm, booktabs, tabularx, array, url, hyperref; the wrapper is the lane's pinned \texttt{amsart} editorial wrapper at 10pt on A4, which restates the theorem-like environments the retained text needs and adds the article's own macro definitions that the retained text needs (\texttt{\textbackslash F}, \texttt{\textbackslash FDP}, \texttt{\textbackslash Rfam}, \texttt{\textbackslash one}), with extra wrapper packages (bm, bbm, dsfont, xspace, cancel, mathtools). The delivered numbering is the unit's own. Assets: no retained span contains a figure environment, a graphics-inclusion command, a PDF-page inclusion, a table or a TikZ picture; no image file is copied into this package, the package declares no asset dependency and no image file is redistributed with it. Licence: version arXiv:2608.09927v1 of this article is distributed under the Creative Commons Attribution 4.0 International licence, as recorded in the arXiv licence metadata for this exact version and shown on the abstract page https://arxiv.org/abs/2608.09927v1. A full rights notice is delivered at licenses/arxiv-2608.09927-CC-BY-4.0.txt. Characters: every retained excerpt is pure ASCII, so the delivered engine, XeLaTeX with its default Latin Modern text font, reports no missing character.
\begin{proposition}[Necessity of coherence for the forward theorem]\label{prop:coherence-necessary}
There exist two true null hypotheses, deterministic active sets, and a dynamic intersection collection for which every fixed-intersection process is a nonnegative test martingale, but the induced dynamic $e$-closure at level $\alpha=1/2$ has stopped FDR equal to one. Hence fixed-intersection $e$-process validity alone does not imply the forward dynamic $e$-closure theorem.
\end{proposition}

\begin{proof}
Let the statistical model consist of one distribution $P$ on $\Omega=\{a,b\}$ with $P(a)=P(b)=1/2$, and let both null hypotheses equal the full model. Put $I_0=\varnothing$, $I_1=\{1\}$, and $I_2=\{1,2\}$, with $I_t=I_2$ thereafter. Let $\F_0$ be trivial and $\F_t=\sigma(\{a\})$ for $t\geq1$.

For each nonempty $S\subseteq\{1,2\}$, set $E_0^S=1$ and, for $t\geq1$, define
\[
 E_t^{\{1\}}=2\one_{\{a\}},\qquad
 E_t^{\{2\}}=2\one_{\{b\}},\qquad
 E_t^{\{1,2\}}=2\one_{\{b\}}.
\]
Each process has conditional expectation one at the first update and is constant afterwards, so every $E^S$ is a nonnegative test martingale for $H_S$.

Define the global stopping time
\[
 \tau=\begin{cases}1,&\omega=a,\\2,&\omega=b.\end{cases}
\]
On $\{a\}$, at time one the set $R=\{1\}$ satisfies its only nonempty closure constraint because
\[
 \FDP_{\{1\}}(R)=1=\tfrac12 E_1^{\{1\}}.
\]
On $\{b\}$, at time two the set $R=\{2\}$ satisfies all three constraints: the $\{1\}$-constraint has loss zero, while
\[
 \FDP_{\{2\}}(R)=\FDP_{\{1,2\}}(R)=1
 =\tfrac12E_2^{\{2\}}=\tfrac12E_2^{\{1,2\}}.
\]
Thus the dynamic closure contains an all-null rejection set at the stopping time on every sample path. Hence
\[
 \max_{R\in\Rfam_{1/2,\tau}(\mathbf E)}\FDP_{\{1,2\}}(R)=1
 \quad\text{almost surely},
\]
and the stopped FDR equals one.

The failure is precisely a coherence violation. At time one on $\{a\}$,
\[
 E_1^{\{1\}}=2>0=E_1^{\{1,2\}},
\]
so the stopping rule selects favourable evidence from different intersection indices on different paths. Each deterministic-index process is valid, but the random-index process is not controlled.
\end{proof}

\end{document}
